
\section{分次环的映射与射影概形的映射}

正如\trRing{}的\trMap{}以相反的方向对应于\trAffineScheme{}的\trMap，（\trBaseRing{} $A$ 上的）\trGradedRing{}的\trMap{}有时也以相反的方向给出\trProjectiveScheme{}的\trMap。这是一次并不完美的推广：并非每个\trGradedRing{}的\trMap{}都给出\trProjectiveScheme{}的\trMap（\S7.4.2）；并非每个\trProjectiveScheme{}的\trMap{}都来自\trGradedRing{}的\trMap（\S19.11.9）；而且不同的\trGradedRing{}的\trMap{}可以给出同一个\trScheme{}的\trMap（练习~7.4.C）。

在做下面这道练习时，把例~7.4.1 与例~7.4.2 想上一遍，或许不无助益。

\textbf{7.4.A.（必作练习）}\quad 设 $\phi\colon S_\bullet\longrightarrow R_\bullet$ 是 ($\mathbb{Z}^{\ge0}$-)\trGradedRing{}之间的一个\trMorphism。（所谓\textbf{\trGradedRing{}映射}（map of graded rings），我们指的是这样一个\trRing{}的\trMap：它作为「\trGrading{}\trAbelianSemigroup」（或「\trGrading{}\trCommutativeMonoid」）之间的\trMap{}保持\trGrading。换言之，存在某个 $d>0$，使得对一切 $n$，$S_n$ 都被映入 $R_{dn}$。）证明：它诱导\trScheme{}的一个\trMorphism
\[
(\Proj R_\bullet)\setminus V(\phi(S_+))\longrightarrow\Proj S_\bullet.
\]
（提示：设 $f$ 是 $S_+$ 的一个\trHomogeneous{}元素。定义一个\trMap{} $D_+(\phi(f))\to D_+(f)$。证明它们可以粘合起来（当 $f$ 跑遍 $S_+$ 的全部\trHomogeneous{}元素时）。证明这就定义了一个从整个 $\Proj R_\bullet\setminus V(\phi(S_+))$ 出发的\trMap。）特别地，若
\begin{equation*}\tag{7.4.0.1}\label{eq:vphi-s-plus-empty}
V(\phi(S_+))=\varnothing,
\end{equation*}
则我们便有\trMorphism{} $\Proj R_\bullet\to\Proj S_\bullet$。从你的解答可以看出：若 $\phi$ 更是 $A$-\trAlgebra{}之间的\trMorphism，那么所诱导的\trMorphism{} $\Proj R_\bullet\setminus V(\phi(S_+))\to\Proj S_\bullet$ 就是 $A$-\trScheme{}之间的\trMorphism。

\textbf{7.4.1.}\ \emph{例。}\quad 让我们看看练习~7.4.A 的实际运用。我们将以概形论的方式解读复射影流形之间的如下\trMap
\[
\mathbb{CP}^1\longrightarrow\mathbb{CP}^2
\]
\[
[s,t]\longmapsto[s^{20},s^9t^{11},t^{20}].
\]
首先注意，这是\trWellDefined{}的：$[\lambda s,\lambda t]$ 与 $[s,t]$ 被送到 $\mathbb{CP}^2$ 中同一个点。它之所以\trWellDefined，缘由在于：$s^{20}$、$s^9t^{11}$、$t^{20}$ 这三个多项式都是次数 $20$ 的\trHomogeneous{}多项式。

在代数上，这对应于一个方向相反的\trGradedRing{}的\trMap
\[
\mathbb{C}[x,y,z]\longrightarrow\mathbb{C}[s,t],
\]
由 $x\mapsto s^{20}$, $y\mapsto s^9t^{11}$, $z\mapsto t^{20}$ 给出。你应当依自己对练习~7.4.A 的解答来解读这件事，并与 \S3.2.11 的仿射例子相比较。

\textbf{7.4.2.}\ \emph{例。}\quad 注意，并不存在由 $[x,y,z]\mapsto[x,y]$ 给出的复\trManifold{}之间的\trMap{} $\mathbb{CP}^2\to\mathbb{CP}^1$，因为当 $x=y=0$ 时该\trMap{}没有定义。这对应于如下事实：由 $s\mapsto x$ 与 $t\mapsto y$ 给出的\trGradedRing{}的\trMap{} $\mathbb{C}[s,t]\to\mathbb{C}[x,y,z]$ 不满足假设 \eqref{eq:vphi-s-plus-empty}。

\textbf{7.4.B.（练习）}\quad 证明：若 $\phi\colon S_\bullet\to R_\bullet$ 满足 $\sqrt{(\phi(S_+))}=R_+$，则假设 \eqref{eq:vphi-s-plus-empty} 成立。（提示：见\cref{ex:empty-vanishing-equivalents}。）这个更为几何化的假设的代数表述，有时验证起来更为容易。

\textbf{7.4.C.（不重要的练习）}\quad 这道练习表明：不同的\trGradedRing{}的\trMap{}可以给出同一个\trScheme{}的\trMap。设 $R_\bullet=k[x,y,z]/(xz,yz,z^2)$、$S_\bullet=k[a,b,c]/(ac,bc,c^2)$，其中每个变量的次数都是 $1$。证明：$\Proj R_\bullet\cong\Proj S_\bullet\cong\mathbb{P}^1_k$。证明：由 $(a,b,c)\mapsto(x,y,z)$ 与 $(a,b,c)\mapsto(x,y,0)$ 给出的那两个\trMap{} $S_\bullet\to R_\bullet$，给出同一个（同）构 $\Proj R_\bullet\to\Proj S_\bullet$。（真正的缘由是：所有这些构造对「在有穷多个次数上发生的事」都不敏感。这一点日后将从多个途径说得精确，最近的便是练习~7.4.F。）

\textbf{7.4.3.}\ \emph{不重要的注记。}\quad \mbox{练习~15.2.H} 表明：即便在相当合理的情形下，也并非\trScheme{}的每个\trMorphism{} $\Proj R_\bullet\to\Proj S_\bullet$ 都来自\trGradedRing{}的\trMap{} $S_\bullet\to R_\bullet$。

\textbf{7.4.4.\ \trVeroneseSubring{}。}\quad 这里有一个有用的构造。设 $S_\bullet$ 是一个\trFinitelyGenerated\trGradedRing。把 $S_\bullet$ 的第 $n$ 个\textbf{\trVeroneseSubring{}}（$n$th Veronese subring）定义为 $S_{n\bullet}=\oplus_{j=0}^{\infty}S_{nj}$。（「旧次数」$n$ 即「新次数」$1$。）这一构造的几何解读，就是重要的\emph{\trVeroneseEmbedding{}}（Veronese embedding），将于 \S9.3.6 中讨论。

\textbf{7.4.D.（练习）}\quad 证明：\trGradedRing{}的\trMap{} $S_{n\bullet}\hookrightarrow S_\bullet$ 诱导一个\emph{\trIsomorphism} $\Proj S_\bullet\mathrel{\xrightarrow{\ \sim\ }}\Proj S_{n\bullet}$。（提示：若 $f\in S_+$ 是次数被 $n$ 整除的\trHomogeneous{}元素，把 $\Proj S_\bullet$ 上的 $D(f)$ 与 $\Proj S_{n\bullet}$ 上的 $D(f)$ 等同起来。为什么这样的\trDistinguishedOpen{}覆盖 $\Proj S_\bullet$？）

\textbf{7.4.E.（练习）}\quad 若 $S_\bullet$ 在次数 $1$ 处生成，证明 $S_{n\bullet}$ 也在次数 $1$ 处生成。（你或许愿意先考虑\trPolynomialRing{}的情形。）

\textbf{7.4.F.（练习）}\quad 证明：若 $R_\bullet$ 与 $S_\bullet$ 是同一个\trFinitelyGenerated\trGradedRing，仅在有穷多个非零次数上彼此不同（请把这话说精确！），则 $\Proj R_\bullet\cong\Proj S_\bullet$。

\textbf{7.4.G.（练习）}\quad 设 $S_\bullet$ 在 $S_0$ 上由 $f_1,\dots,f_n$ 生成。找出一个 $d$，使得 $S_{d\bullet}$ 在「新」次数 $1$（即「旧」次数 $d$）处是\trFinitelyGenerated{}的。（这出人意料地棘手，所以给一条提示。设有\trGenerator{} $x_1,\dots,x_n$，次数依次为 $d_1,\dots,d_n$。证明：任何次数至少为 $nd_1\cdots d_n$ 的\trMonomial{} $x_1^{\alpha_1}\cdots x_n^{\alpha_n}$，必有某个 $i$ 使得 $\alpha_i\ge(\prod_j d_j)/d_i$。再证明：第 $nd_1\cdots d_n$ 个\trVeroneseSubring{}由「新」次数 $1$ 处的元素生成。）

练习~7.4.G 与练习~7.4.D 结合起来表明：假定\trFinitelyGenerated\trGradedRing{}都在次数 $1$ 处生成，几乎无害——因为在\trRegrading{}之后（更精确地说：只保留次数为 $d$ 之倍数的那些项，再把次数除以 $d$），情形确实如此。这一点很便利，因为它意味着：借练习~7.4.D 之助，我们可以假定任何\trFinitelyGenerated\trGradedRing{}都在次数 $1$ 处生成。\mbox{练习~9.3.F} 日后将作为一个推论蕴含：我们可以把每个 $\Proj$ 嵌入某个\trProjectiveSpace{}之中。

\textbf{7.4.H.（较不重要的练习）}\quad 设 $S_\bullet$ 是一个\trFinitelyGenerated\trGradedRing。证明：$S_{n\bullet}$ 是一个\trFinitelyGenerated\trGradedRing。（一个可能的途径：用上一道练习，或某个类似的办法，证明存在某个 $N$ 使得 $S_{nN\bullet}$ 在次数 $1$ 处生成，从而\trGradedRing{} $S_{nN\bullet}$ 是\trFinitelyGenerated{}的。然后证明：对每个 $0<j<N$，$S_{nN\bullet+nj}$ 是 $S_{nN\bullet}$ 上的一个\trFinitelyGenerated{}的\trModule。）
